\chapter{Аналитический раздел}

\section{Постановка задачи}

JPEG-сжатие с потерями изменяет значения пикселей изображения. Если полученный файл декодировать и снова сохранить в JPEG, изображение пройдёт второй цикл сжатия. Следы этих операций могут использоваться при анализе истории обработки изображения. Повторное сохранение возможно и без изменения содержания снимка, поэтому обнаружение двукратного сжатия само по себе не доказывает наличие подделки.

В работе решается задача классификации JPEG-файла по числу выполненных сжатий. Рассматриваются два класса:
\begin{itemize}
  \item однократно сжатые изображения, полученные JPEG-кодированием исходных изображений без предшествующего JPEG-сжатия;
  \item двукратно сжатые изображения, полученные декодированием однократно сжатых изображений и их повторным JPEG-кодированием.
\end{itemize}

Входными данными служит только проверяемый JPEG-файл. Исходное несжатое изображение и параметры первого сжатия при проверке неизвестны. Выходными данными являются предсказанный класс и оценка модели для класса двукратного сжатия.

Область исследования ограничивается повторным сжатием с совпадающей сеткой блоков JPEG. Между сохранениями не выполняются масштабирование, поворот и смещение изображения. Предполагается, что всё изображение имеет общую историю сжатия. Локализация вставленных областей и определение точного числа сжатий, превышающего два, в задачу не входят.

\section{Признаки повторного JPEG-сжатия}

В режиме JPEG на основе дискретного косинусного преобразования (ДКП) компоненты изображения обрабатываются блоками размером $8\times8$ отсчётов. Каждый блок преобразуется в набор частотных коэффициентов, которые затем квантуются и кодируются~\cite{jpeg_standard}. Квантование уменьшает точность представления коэффициентов и создаёт основные потери информации.

Для отдельного коэффициента операция квантования записывается в виде
\begin{equation}
  k = \operatorname{round}\left(\frac{c}{q}\right),
\end{equation}
где $c$ --- коэффициент ДКП до квантования; $q$ --- соответствующий элемент таблицы квантования; $k$ --- целочисленный результат; $\operatorname{round}$ --- округление до ближайшего целого. При декодировании восстанавливается значение $kq$, которое в общем случае отличается от $c$.

После второго сжатия распределения коэффициентов могут приобретать дополнительные максимумы и пропуски значений. Эти закономерности зависят от таблиц первого и второго квантования. Для их описания используют гистограммы: для каждой выбранной частоты подсчитывается, как часто коэффициенты попадают в заданные интервалы~\cite{wang2016}. Нейронная сеть может обучаться различать такие распределения у двух классов изображений.

В JPEG-файле доступны таблицы квантования последнего сохранения. Они характеризуют условия получения наблюдаемого изображения и могут передаваться модели вместе с гистограммами~\cite{park2018}. Числовой параметр качества, задаваемый в графическом редакторе, служит способом выбора таблиц и не является универсальной характеристикой для всех программ кодирования.

Выраженность следов повторного сжатия неодинакова. При близких параметрах двух сохранений различия могут быть малы. Однородные области также содержат меньше частотной информации для классификации~\cite{park2018}. Поэтому работоспособность метода определяется экспериментально на размеченных данных, а результат для отдельного файла может содержать ошибку.

\section{Обзор существующих решений}

Рассматриваются два нейросетевых метода, различающихся составом частотных признаков и способом учёта параметров JPEG-сжатия.

\subsection{Метод Wang--Zhang}

Q. Wang и R. Zhang предложили свёрточную нейронную сеть, входом которой служат гистограммы низкочастотных коэффициентов ДКП~\cite{wang2016}. Признаки объединяются в вектор из 99 значений. Модель определяет, был ли анализируемый фрагмент сжат один или два раза.

В опубликованном эксперименте обучались отдельные классификаторы для разных значений параметра качества последнего сжатия. Такой подход позволяет учитывать заранее выбранные условия кодирования, но требует набора моделей для нескольких значений качества. Работа также рассматривает поиск областей с различной историей сжатия посредством последовательной проверки фрагментов изображения.

Для разрабатываемой программы из этого метода полезен принцип компактного описания изображения через гистограммы. Поддержка отдельных моделей для каждого значения качества усложнила бы выбор классификатора при проверке произвольного входного файла.

\subsection{Метод Park и реализация DJPEG-torch}

J. Park и соавторы предложили совместно использовать гистограммы коэффициентов ДКП и таблицу квантования из заголовка JPEG~\cite{park2018}. Модель обучается на изображениях со смешанными параметрами сжатия. В опубликованной архитектуре предусмотрены четыре свёрточных слоя и три полносвязных слоя.

Авторы предоставляют реализацию DJPEG-torch, содержащую обучение, проверку модели и пример поиска изменённых областей~\cite{djpeg_code}. Входные изображения преобразуются в яркостную компоненту; по блокам этой компоненты вычисляются ДКП и гистограммы. Таблица квантования дополнительно поступает в полносвязные слои.

Преимущество решения для рассматриваемой задачи состоит в учёте параметров последнего сохранения внутри одной модели. Реализация пригодна для изучения и повторного обучения. Однако её использование само по себе не отвечает на вопрос, насколько можно упростить классификатор при сохранении приемлемого качества. Для этого требуется сравнение с компактными архитектурами на общих данных.

\section{Сравнение решений и выбор подхода}

Критериями сравнения выбраны входные признаки, способ учёта параметров сжатия и пригодность решения для проведения собственного эксперимента. Сводные характеристики приведены в таблице~\ref{tab:analogues}.

\begin{table}[H]
\caption{Сравнение нейросетевых методов}
\label{tab:analogues}
\small\singlespacing
\renewcommand{\arraystretch}{1.25}
\begin{tabularx}{\textwidth}{|>{\raggedright\arraybackslash}p{0.24\textwidth}|>{\raggedright\arraybackslash}X|>{\raggedright\arraybackslash}X|}
\hline
Критерий & Wang--Zhang~\cite{wang2016} & Park / DJPEG-torch~\cite{park2018,djpeg_code} \\ \hline
Входные признаки & Гистограммы выбранных коэффициентов ДКП & Гистограммы ДКП и таблица квантования \\ \hline
Учёт последнего сжатия & Отдельные классификаторы для выбранных значений качества & Таблица квантования подаётся одной модели \\ \hline
Основа для собственной реализации & Компактные частотные признаки & Совместная обработка признаков и параметров сжатия; опубликованный код \\ \hline
Что необходимо проверить в работе & Пригодность компактного классификатора для смешанных параметров & Соотношение качества и времени у более простых моделей \\ \hline
\end{tabularx}
\end{table}

Сравнивать численные оценки качества из разных публикаций напрямую некорректно: они получены при различных параметрах сжатия и составе выборок. Время обработки также зависит от оборудования и учитываемых этапов вычислений. Поэтому в работе модели будут обучены на одной выборке и проверены по единому протоколу.

Предлагается разработать программу с общим извлечением признаков и двумя заменяемыми классификаторами: многослойным персептроном и одномерной свёрточной нейронной сетью. Персептрон рассматривается как базовый вариант, а свёрточная сеть --- как способ учитывать соседство интервалов гистограмм. Обе модели будут получать одинаковую информацию о JPEG-файле, включая параметры последнего квантования.

Собственная реализация нужна для самостоятельного обучения выбранных моделей, единой проверки изображений и сопоставимых измерений. Существующие методы дают основу для выбора признаков. Преимущество одной из исследуемых архитектур должно определяться по результатам эксперимента.

\section{Выбор данных и требований к программе}

Для обучения выбран опубликованный набор данных авторов DJPEG-torch, содержащий однократно и двукратно сжатые фрагменты размером $256\times256$ пикселей~\cite{djpeg_data,djpeg_dataset}. Он позволяет использовать готовую разметку без ручного определения истории каждого файла. В работе будет сформирована ограниченная подвыборка с равным числом объектов двух классов.

При подготовке данных необходимо исключить пересечение одинаковых фрагментов и их версий между обучающей, валидационной и тестовой выборками. Возможность объединения разных фрагментов одной фотографии определяется доступными метаданными набора. Если такую связь установить нельзя, результаты будут характеризовать классификацию фрагментов данного набора, а перенос на новые фотографии потребует отдельной проверки.

Авторы набора указывают, что при его формировании изменяли квантование яркостной компоненты, не выполняя аналогичное квантование цветоразностных компонент~\cite{djpeg_data}. Поэтому в работе используются яркостные признаки. Эта особенность ограничивает перенос полученных оценок на JPEG-файлы, созданные другими способами.

К программе предъявляются следующие требования:
\begin{enumerate}
  \item принимать JPEG-файл и выполнять проверку без доступа к несжатому оригиналу;
  \item обрабатывать изображения размером не менее $256\times256$ пикселей; для больших изображений выделять фрагмент этого размера с сохранением сетки блоков JPEG;
  \item позволять выбрать одну из двух самостоятельно обученных нейронных сетей;
  \item выдавать класс изображения и оценку модели для двукратного сжатия;
  \item сохранять результаты проверки тестовой выборки и измерять полное время обработки файла, включая подготовку признаков;
  \item сообщать об ошибке для повреждённых файлов и неподдерживаемых входных данных.
\end{enumerate}

Качество моделей будет сравниваться по доле правильных ответов и F1-мере для класса двукратного сжатия. Матрица ошибок позволит отдельно оценить ложные обнаружения и пропуски. Быстродействие будет измеряться на одном устройстве при последовательной обработке изображений.

\section*{Выводы}
\addcontentsline{toc}{section}{Выводы}

Для исследования выбраны гистограммы коэффициентов ДКП с дополнительной информацией о таблице квантования. На этих признаках планируется обучить персептрон и одномерную свёрточную сеть. Готовый размеченный набор DJPEG-torch позволяет провести сравнение при общей подготовке данных.
